
\documentclass{article}
\usepackage{colortbl}
\usepackage{makecell}
\usepackage{multirow}
\usepackage{supertabular}

\begin{document}

\newcounter{utterance}

\centering \large Interaction Transcript for game `dond', experiment `coop\_en', episode 11 with kimi{-}k2.5{-}without{-}reasoning{-}t1.0.
\vspace{24pt}

{ \footnotesize  \setcounter{utterance}{1}
\setlength{\tabcolsep}{0pt}
\begin{supertabular}{c@{$\;$}|p{.15\linewidth}@{}p{.15\linewidth}p{.15\linewidth}p{.15\linewidth}p{.15\linewidth}p{.15\linewidth}}
    \# & $\;$A & \multicolumn{4}{c}{Game Master} & $\;\:$B\\
    \hline

    \theutterance \stepcounter{utterance}  
    & & \multicolumn{4}{p{0.6\linewidth}}{
        \cellcolor[rgb]{0.9,0.9,0.9}{
            \makecell[{{p{\linewidth}}}]{
                \texttt{\tiny{[P1$\langle$GM]}}
                \texttt{You are playing a cooperative negotiation game in which you have to agree on how to divide a set of items among you and another player.} \\
\\ 
\texttt{Rules:} \\
\texttt{(a) You and the other player are given a set of items. Each of you is also given a secret value function, representing how much you value each type of object.} \\
\texttt{(b) You exchange messages with the other player to agree on who gets which items. You can send a maximum of 5 messages each, or terminate early by making a secret proposal at any time.} \\
\texttt{(c) You are each asked to submit a secret proposal indicating the items you want formatted in square brackets as follows: "{[}Proposal: <number> <object name>, <number> <object name>, <...>{]}"} \\
\texttt{(d) If your proposals are complementary, i.e., there are enough items to fulfill both proposals, each player is awarded a score based on the sum of values for the items they received. Otherwise, both of you get zero points.} \\
\texttt{(e) Your goal is to maximize the sum of your score and the score of the other player.} \\
\\ 
\texttt{Let us start.} \\
\\ 
\texttt{The set of available items is:} \\
\\ 
\texttt{1 ring, 3 fruits, 1 chalk} \\
\\ 
\texttt{Your secret value function is:} \\
\\ 
\texttt{ring: 6, fruit: 1, chalk: 1} \\
\\ 
\texttt{IMPORTANT: Your messages, unless it is the secret proposal, are directly transmitted to the other player, so do not include any response to the rules or text announcing your message. To make a secret proposal, use the indicated format. Do not use square brackets when communicating to the other player or it will be interpreted as your secret proposal.} \\
            }
        }
    }
    & \\ \\

    \theutterance \stepcounter{utterance}  
    & \multicolumn{4}{p{0.6\linewidth}}{
        \cellcolor[rgb]{0.9,0.9,0.9}{
            \makecell[{{p{\linewidth}}}]{
                \texttt{\tiny{[P1$\rangle$GM]}}
                \texttt{I value the ring highly (6) and fruits and chalk equally (1 each). Since I want to maximize total score and the other player likely has different values, I should try to find out what they want.} \\
\\ 
\texttt{Let me start by asking about their preferences to find a mutually beneficial split.} \\
\\ 
\texttt{What items do you value most? I really want the ring but am flexible on the fruits and chalk.} \\
            }
        }
    }
    & & \\ \\

    \theutterance \stepcounter{utterance}  
    & & & \multicolumn{2}{p{0.3\linewidth}}{
        \cellcolor[rgb]{0.9,0.9,0.9}{
            \makecell[{{p{\linewidth}}}]{
                \texttt{\tiny{[GM$|$GM]}}
                \texttt{continue} \\
            }
        }
    }
    & & \\ \\

    \theutterance \stepcounter{utterance}  
    & & & \multicolumn{2}{p{0.3\linewidth}}{
        \cellcolor[rgb]{0.9,0.9,0.9}{
            \makecell[{{p{\linewidth}}}]{
                \texttt{\tiny{[GM$|$GM]}}
                \texttt{I value the ring highly (6) and fruits and chalk equally (1 each). Since I want to maximize total score and the other player likely has different values, I should try to find out what they want.} \\
\\ 
\texttt{Let me start by asking about their preferences to find a mutually beneficial split.} \\
\\ 
\texttt{What items do you value most? I really want the ring but am flexible on the fruits and chalk.} \\
            }
        }
    }
    & & \\ \\

    \theutterance \stepcounter{utterance}  
    & & \multicolumn{4}{p{0.6\linewidth}}{
        \cellcolor[rgb]{0.9,0.9,0.9}{
            \makecell[{{p{\linewidth}}}]{
                \texttt{\tiny{[GM$\rangle$P2]}}
                \texttt{You are playing a cooperative negotiation game in which you have to agree on how to divide a set of items among you and another player.} \\
\\ 
\texttt{Rules:} \\
\texttt{(a) You and the other player are given a set of items. Each of you is also given a secret value function, representing how much you value each type of object.} \\
\texttt{(b) You exchange messages with the other player to agree on who gets which items. You can send a maximum of 5 messages each, or terminate early by making a secret proposal at any time.} \\
\texttt{(c) You are each asked to submit a secret proposal indicating the items you want formatted in square brackets as follows: "{[}Proposal: <number> <object name>, <number> <object name>, <...>{]}"} \\
\texttt{(d) If your proposals are complementary, i.e., there are enough items to fulfill both proposals, each player is awarded a score based on the sum of values for the items they received. Otherwise, both of you get zero points.} \\
\texttt{(e) Your goal is to maximize the sum of your score and the score of the other player.} \\
\\ 
\texttt{Let us start.} \\
\\ 
\texttt{The set of available items is:} \\
\\ 
\texttt{1 ring, 3 fruits, 1 chalk} \\
\\ 
\texttt{Your secret value function is:} \\
\\ 
\texttt{ring: 6, fruit: 1, chalk: 1} \\
\\ 
\texttt{IMPORTANT: Your messages, unless it is the secret proposal, are directly transmitted to the other player, so do not include any response to the rules or text announcing your message. To make a secret proposal, use the indicated format. Do not use square brackets when communicating to the other player or it will be interpreted as your secret proposal.} \\
\\ 
\\ 
\texttt{I value the ring highly (6) and fruits and chalk equally (1 each). Since I want to maximize total score and the other player likely has different values, I should try to find out what they want.} \\
\\ 
\texttt{Let me start by asking about their preferences to find a mutually beneficial split.} \\
\\ 
\texttt{What items do you value most? I really want the ring but am flexible on the fruits and chalk.} \\
            }
        }
    }
    & \\ \\

    \theutterance \stepcounter{utterance}  
    & & & \multicolumn{4}{p{0.6\linewidth}}{
        \cellcolor[rgb]{0.9,0.9,0.9}{
            \makecell[{{p{\linewidth}}}]{
                \texttt{\tiny{[GM$\langle$P2]}}
                \texttt{I value the ring highly and am flexible on the fruits and chalk. What do you prefer?} \\
            }
        }
    }
     \\ \\

    \theutterance \stepcounter{utterance}  
    & & & \multicolumn{2}{p{0.3\linewidth}}{
        \cellcolor[rgb]{0.9,0.9,0.9}{
            \makecell[{{p{\linewidth}}}]{
                \texttt{\tiny{[GM$|$GM]}}
                \texttt{continue} \\
            }
        }
    }
    & & \\ \\

    \theutterance \stepcounter{utterance}  
    & & & \multicolumn{2}{p{0.3\linewidth}}{
        \cellcolor[rgb]{0.9,0.9,0.9}{
            \makecell[{{p{\linewidth}}}]{
                \texttt{\tiny{[GM$|$GM]}}
                \texttt{I value the ring highly and am flexible on the fruits and chalk. What do you prefer?} \\
            }
        }
    }
    & & \\ \\

    \theutterance \stepcounter{utterance}  
    & & \multicolumn{4}{p{0.6\linewidth}}{
        \cellcolor[rgb]{0.9,0.9,0.9}{
            \makecell[{{p{\linewidth}}}]{
                \texttt{\tiny{[P1$\langle$GM]}}
                \texttt{I value the ring highly and am flexible on the fruits and chalk. What do you prefer?} \\
            }
        }
    }
    & \\ \\

    \theutterance \stepcounter{utterance}  
    & \multicolumn{4}{p{0.6\linewidth}}{
        \cellcolor[rgb]{0.9,0.9,0.9}{
            \makecell[{{p{\linewidth}}}]{
                \texttt{\tiny{[P1$\rangle$GM]}}
                \texttt{We both value the ring highly, which is tricky. Since I value fruits and chalk equally at 1 each, and you said you're flexible on those, would you consider letting me have the ring if I give you all 3 fruits and the chalk? That gives you 4 items while I get 1, but the ring is worth a lot to me.} \\
\\ 
\texttt{Alternatively, we could split: I take the ring and 1 fruit, you take 2 fruits and the chalk?} \\
            }
        }
    }
    & & \\ \\

    \theutterance \stepcounter{utterance}  
    & & & \multicolumn{2}{p{0.3\linewidth}}{
        \cellcolor[rgb]{0.9,0.9,0.9}{
            \makecell[{{p{\linewidth}}}]{
                \texttt{\tiny{[GM$|$GM]}}
                \texttt{continue} \\
            }
        }
    }
    & & \\ \\

    \theutterance \stepcounter{utterance}  
    & & & \multicolumn{2}{p{0.3\linewidth}}{
        \cellcolor[rgb]{0.9,0.9,0.9}{
            \makecell[{{p{\linewidth}}}]{
                \texttt{\tiny{[GM$|$GM]}}
                \texttt{We both value the ring highly, which is tricky. Since I value fruits and chalk equally at 1 each, and you said you're flexible on those, would you consider letting me have the ring if I give you all 3 fruits and the chalk? That gives you 4 items while I get 1, but the ring is worth a lot to me.} \\
\\ 
\texttt{Alternatively, we could split: I take the ring and 1 fruit, you take 2 fruits and the chalk?} \\
            }
        }
    }
    & & \\ \\

    \theutterance \stepcounter{utterance}  
    & & \multicolumn{4}{p{0.6\linewidth}}{
        \cellcolor[rgb]{0.9,0.9,0.9}{
            \makecell[{{p{\linewidth}}}]{
                \texttt{\tiny{[GM$\rangle$P2]}}
                \texttt{We both value the ring highly, which is tricky. Since I value fruits and chalk equally at 1 each, and you said you're flexible on those, would you consider letting me have the ring if I give you all 3 fruits and the chalk? That gives you 4 items while I get 1, but the ring is worth a lot to me.} \\
\\ 
\texttt{Alternatively, we could split: I take the ring and 1 fruit, you take 2 fruits and the chalk?} \\
            }
        }
    }
    & \\ \\

    \theutterance \stepcounter{utterance}  
    & & & \multicolumn{4}{p{0.6\linewidth}}{
        \cellcolor[rgb]{0.9,0.9,0.9}{
            \makecell[{{p{\linewidth}}}]{
                \texttt{\tiny{[GM$\langle$P2]}}
                \texttt{I appreciate the offer, but I also value the ring very highly. Let me propose a different split: what if you take all 3 fruits and the chalk (total value 4 for you based on your stated values), and I take just the ring? That gives you 4 points and me 6 points, for a total of 10.} \\
\\ 
\texttt{Or if you value the ring even higher than 6, we could discuss a different arrangement. What are your actual values for each item?} \\
            }
        }
    }
     \\ \\

    \theutterance \stepcounter{utterance}  
    & & & \multicolumn{2}{p{0.3\linewidth}}{
        \cellcolor[rgb]{0.9,0.9,0.9}{
            \makecell[{{p{\linewidth}}}]{
                \texttt{\tiny{[GM$|$GM]}}
                \texttt{continue} \\
            }
        }
    }
    & & \\ \\

    \theutterance \stepcounter{utterance}  
    & & & \multicolumn{2}{p{0.3\linewidth}}{
        \cellcolor[rgb]{0.9,0.9,0.9}{
            \makecell[{{p{\linewidth}}}]{
                \texttt{\tiny{[GM$|$GM]}}
                \texttt{I appreciate the offer, but I also value the ring very highly. Let me propose a different split: what if you take all 3 fruits and the chalk (total value 4 for you based on your stated values), and I take just the ring? That gives you 4 points and me 6 points, for a total of 10.} \\
\\ 
\texttt{Or if you value the ring even higher than 6, we could discuss a different arrangement. What are your actual values for each item?} \\
            }
        }
    }
    & & \\ \\

    \theutterance \stepcounter{utterance}  
    & & \multicolumn{4}{p{0.6\linewidth}}{
        \cellcolor[rgb]{0.9,0.9,0.9}{
            \makecell[{{p{\linewidth}}}]{
                \texttt{\tiny{[P1$\langle$GM]}}
                \texttt{I appreciate the offer, but I also value the ring very highly. Let me propose a different split: what if you take all 3 fruits and the chalk (total value 4 for you based on your stated values), and I take just the ring? That gives you 4 points and me 6 points, for a total of 10.} \\
\\ 
\texttt{Or if you value the ring even higher than 6, we could discuss a different arrangement. What are your actual values for each item?} \\
            }
        }
    }
    & \\ \\

    \theutterance \stepcounter{utterance}  
    & \multicolumn{4}{p{0.6\linewidth}}{
        \cellcolor[rgb]{0.9,0.9,0.9}{
            \makecell[{{p{\linewidth}}}]{
                \texttt{\tiny{[P1$\rangle$GM]}}
                \texttt{I see you're trying to get me to reveal my exact values, but I already stated them: ring 6, fruit 1, chalk 1. Your proposal gives me 4 points (3 fruits + chalk) and you 6 points (ring), total 10.} \\
\\ 
\texttt{But wait {-} if I took the ring (6) and gave you all 3 fruits and chalk (4), that's also total 10. Same outcome.} \\
\\ 
\texttt{Since we both value the ring at 6, we need to decide who gets it. Given that we value it equally, and I proposed first to give you all 4 other items, would you accept: you get 3 fruits + chalk, I get ring? Or do you have a reason you should get the ring instead?} \\
            }
        }
    }
    & & \\ \\

    \theutterance \stepcounter{utterance}  
    & & & \multicolumn{2}{p{0.3\linewidth}}{
        \cellcolor[rgb]{0.9,0.9,0.9}{
            \makecell[{{p{\linewidth}}}]{
                \texttt{\tiny{[GM$|$GM]}}
                \texttt{continue} \\
            }
        }
    }
    & & \\ \\

    \theutterance \stepcounter{utterance}  
    & & & \multicolumn{2}{p{0.3\linewidth}}{
        \cellcolor[rgb]{0.9,0.9,0.9}{
            \makecell[{{p{\linewidth}}}]{
                \texttt{\tiny{[GM$|$GM]}}
                \texttt{I see you're trying to get me to reveal my exact values, but I already stated them: ring 6, fruit 1, chalk 1. Your proposal gives me 4 points (3 fruits + chalk) and you 6 points (ring), total 10.} \\
\\ 
\texttt{But wait {-} if I took the ring (6) and gave you all 3 fruits and chalk (4), that's also total 10. Same outcome.} \\
\\ 
\texttt{Since we both value the ring at 6, we need to decide who gets it. Given that we value it equally, and I proposed first to give you all 4 other items, would you accept: you get 3 fruits + chalk, I get ring? Or do you have a reason you should get the ring instead?} \\
            }
        }
    }
    & & \\ \\

    \theutterance \stepcounter{utterance}  
    & & \multicolumn{4}{p{0.6\linewidth}}{
        \cellcolor[rgb]{0.9,0.9,0.9}{
            \makecell[{{p{\linewidth}}}]{
                \texttt{\tiny{[GM$\rangle$P2]}}
                \texttt{I see you're trying to get me to reveal my exact values, but I already stated them: ring 6, fruit 1, chalk 1. Your proposal gives me 4 points (3 fruits + chalk) and you 6 points (ring), total 10.} \\
\\ 
\texttt{But wait {-} if I took the ring (6) and gave you all 3 fruits and chalk (4), that's also total 10. Same outcome.} \\
\\ 
\texttt{Since we both value the ring at 6, we need to decide who gets it. Given that we value it equally, and I proposed first to give you all 4 other items, would you accept: you get 3 fruits + chalk, I get ring? Or do you have a reason you should get the ring instead?} \\
            }
        }
    }
    & \\ \\

    \theutterance \stepcounter{utterance}  
    & & & \multicolumn{4}{p{0.6\linewidth}}{
        \cellcolor[rgb]{0.9,0.9,0.9}{
            \makecell[{{p{\linewidth}}}]{
                \texttt{\tiny{[GM$\langle$P2]}}
                \texttt{You're right that both splits give total 10. Since we value the ring equally, fairness suggests we could flip a coin or find another way to decide. But let me check: is there any split where we both do better?} \\
\\ 
\texttt{Actually, no {-} the ring is the only high{-}value item and we both value it at 6. One of us must get it.} \\
\\ 
\texttt{Since you proposed first and the total is the same either way, I'm willing to accept your proposal: you get the ring, I get 3 fruits and chalk. This gives us 10 total points and is fair given our equal valuations.} \\
\\ 
\texttt{{[}Proposal: 3 fruit, 1 chalk{]}} \\
            }
        }
    }
     \\ \\

    \theutterance \stepcounter{utterance}  
    & & & \multicolumn{2}{p{0.3\linewidth}}{
        \cellcolor[rgb]{0.9,0.9,0.9}{
            \makecell[{{p{\linewidth}}}]{
                \texttt{\tiny{[GM$|$GM]}}
                \texttt{proposal} \\
            }
        }
    }
    & & \\ \\

    \theutterance \stepcounter{utterance}  
    & & & \multicolumn{2}{p{0.3\linewidth}}{
        \cellcolor[rgb]{0.9,0.9,0.9}{
            \makecell[{{p{\linewidth}}}]{
                \texttt{\tiny{[GM$|$GM]}}
                [0, 3, 1]
            }
        }
    }
    & & \\ \\

    \theutterance \stepcounter{utterance}  
    & & \multicolumn{4}{p{0.6\linewidth}}{
        \cellcolor[rgb]{0.9,0.9,0.9}{
            \makecell[{{p{\linewidth}}}]{
                \texttt{\tiny{[P1$\langle$GM]}}
                \texttt{The other player has submitted a secret proposal. It is now your turn to submit a proposal of your own.} \\
            }
        }
    }
    & \\ \\

    \theutterance \stepcounter{utterance}  
    & \multicolumn{4}{p{0.6\linewidth}}{
        \cellcolor[rgb]{0.9,0.9,0.9}{
            \makecell[{{p{\linewidth}}}]{
                \texttt{\tiny{[P1$\rangle$GM]}}
                \texttt{{[}Proposal: 1 ring, 0 fruit, 0 chalk{]}} \\
            }
        }
    }
    & & \\ \\

    \theutterance \stepcounter{utterance}  
    & & & \multicolumn{2}{p{0.3\linewidth}}{
        \cellcolor[rgb]{0.9,0.9,0.9}{
            \makecell[{{p{\linewidth}}}]{
                \texttt{\tiny{[GM$|$GM]}}
                \texttt{proposal} \\
            }
        }
    }
    & & \\ \\

    \theutterance \stepcounter{utterance}  
    & & & \multicolumn{2}{p{0.3\linewidth}}{
        \cellcolor[rgb]{0.9,0.9,0.9}{
            \makecell[{{p{\linewidth}}}]{
                \texttt{\tiny{[GM$|$GM]}}
                [1, 0, 0]
            }
        }
    }
    & & \\ \\

    \theutterance \stepcounter{utterance}  
    & & & \multicolumn{2}{p{0.3\linewidth}}{
        \cellcolor[rgb]{0.9,0.9,0.9}{
            \makecell[{{p{\linewidth}}}]{
                \texttt{\tiny{[GM$|$GM]}}
                [[1, 0, 0], [0, 3, 1]]
            }
        }
    }
    & & \\ \\

\end{supertabular}
}

\end{document}
